\vspace{50mm}
\chapter*{Conclusion Générale}
\addcontentsline{toc}{chapter}{Conclusion Générale}
Ce projet de fin d’études, réalisé au sein de l’entreprise Bee Coders, a abouti au développement d’une plateforme web nommée HrNova, dédiée à la digitalisation des processus RH. Ce prototype opérationnel intègre les fonctionnalités essentielles telles que la gestion des offres d’emploi, des candidatures avec un système intelligent de filtrage de CV, la planification d’entretiens, la gestion des congés, le suivi de projet, l’organisation des tâches et la gestion des comptes utilisateurs selon différents rôles. \\
 
Développée selon une approche agile avec la méthode Scrum, cette solution nous a permis de mettre en œuvre nos compétences techniques, méthodologiques et collaboratives dans un contexte professionnel réel. Toutefois, en tant que prototype destiné à une seule entreprise, HrNova présente certaines limites, notamment l'absence d’une architecture multi-tenant et d’un module de communication interne. \\

Parmi les pistes d’amélioration envisagées, il serait pertinent d’ajouter un super administrateur capable de gérer plusieurs entreprises sur la même plateforme. Chaque entreprise pourrait être identifiée par son domaine de messagerie professionnel (par exemple, @ooredoo.com, @vermeg.com), ce qui permettrait de créer une solution RH centralisée, adaptée à un usage à plus grande échelle. D'autres évolutions possibles incluent l’ajout d’un système de messagerie instantanée interne entre les membres d’une équipe et leur manager pour favoriser la collaboration, ainsi que la mise en place d’un envoi automatique de mails aux candidats lors de la planification d’un entretien, renforçant ainsi la réactivité et la fluidité du processus de recrutement. \\



\markboth{Conclusion Générale}{}